\clearpage
\phantomsection% 
\addcontentsline{toc}{chapter}{KATA PENGANTAR}
%\thispagestyle{fancy}

\begin{justifying}
	\large \bfseries \centering \MakeUppercase{Kata Pengantar}\linebreak
	
	\normalsize \normalfont \justifying

	Puji syukur kehadirat Allah SWT/Tuhan Yang Maha Esa atas limpahan rahmat, karunia, serta petunjuk-Nya sehingga penyusunan tugas akhir yang berjudul Rancang Bangun Teleskop Robotik \textit{Alt-Azimuth} dengan Optimasi Estimasi Sudut Menggunakan \textit{Mahony Filter} (Studi Kasus OAIL ITERA ini dapat terselesaikan dengan baik. Tugas akhir ini disusun sebagai salah satu syarat untuk memperoleh gelar Sarjana Teknik pada Program Studi Teknik Informatika, Fakultas Teknologi Industri, Institut Teknologi Sumatera. \par
	
	Penulis menyadari sepenuhnya bahwa penyelesaian tugas akhir ini tidak lepas dari bimbingan, arahan, bantuan, dan dukungan dari berbagai pihak. Oleh karena itu, pada kesempatan ini penulis mengucapkan terima kasih yang sebesar-besarnya kepada: \par
	\begin{enumerate}
		\item {Prof. Dr. I Nyoman Pugeg Aryantha} selaku Rektor Institut Teknologi Sumatera.  
		\item {Hadi Teguh Yudistira, S.T., Ph.D.} selaku Dekan Fakultas Teknologi Industri.
		\item {Andika Setiawan, S.Kom., M.Cs.} selaku Ketua Program Studi Teknik Informatika.
		\item {Ilham Firman Ashari, S.Kom., M.T.} selaku Koordinator Tugas Akhir Program Studi Teknik Informatika.
		\item \printnamapengujia dan \printnamapengujib selaku Dosen Penguji yang telah memberikan saran, kritik, dan masukan konstruktif untuk penyempurnaan tugas akhir ini.
		\item Seluruh Dosen dan Tenaga Kependidikan Program Studi Teknik Informatika yang telah memberikan ilmu dan pengalaman berharga selama penulis menempuh pendidikan di Institut Teknologi Sumatera.
		\item Kedua orang tua dan seluruh keluarga tercinta yang senantiasa memberikan doa, dukungan moril dan materil, serta motivasi yang tiada henti kepada penulis dalam menyelesaikan studi dan tugas akhir ini.
		\item Semua pihak yang tidak dapat disebutkan satu per satu yang telah membantu penulis dalam menyelesaikan tugas akhir ini.
	\end{enumerate} \par
	
	Penulis menyadari bahwa tugas akhir ini masih jauh dari sempurna dan tidak luput dari kekurangan serta kelemahan. Oleh karena itu, penulis sangat mengharapkan kritik dan saran yang membangun untuk perbaikan di masa mendatang. Akhir kata, penulis berharap semoga tugas akhir ini dapat memberikan manfaat dan kontribusi bagi pengembangan ilmu pengetahuan, serta dapat menjadi referensi bagi penelitian selanjutnya.
	\vfill
	
\end{justifying}
\clearpage





\begin{comment}
\clearpage
\pagestyle{fancy}
\fancyhf{}
\fancyhead[R]{\thepage}
\phantomsection% 
%\clearpage
%\phantomsection% 
\addcontentsline{toc}{chapter}{KATA PENGANTAR}
%\thispagestyle{fancy}

\begin{justifying}
	\large \bfseries \centering \MakeUppercase{Kata Pengantar}\linebreak
	
	\normalsize \normalfont \justifying
	Puji syukur kehadirat Tuhan Yang Maha Esa atas limpahan rahmat, karunia, serta petunjuk-Nya sehingga penyusunan tugas akhir ini telah terselesaikan dengan baik. Dalam penyusunan tugas akhir ini penulis telah banyak mendapatkan arahan, bantuan, serta dukungan dari berbagai pihak. Oleh karena itu pada kesempatan ini penulis mengucapan terima kasih kepada: \par
	\begin{enumerate}
		\item Bapak Prof. Dr. I. Nyoman Pugeg Aryantha selaku Rektor Institut Teknologi Sumatera.  
		\item Bapak Hadi Teguh Yudistira, S.T., Ph.D. selaku Dekan Fakultas Teknologi Industri.
		\item Bapak Andika Setiawan, S. Kom., M. Cs. selaku Ketua Program Studi Teknik Informatika.
		\item Bapak Ilham Firman Ashari, S. Kom., M.T. selaku Koordinator Tugas Akhir Program Studi Teknik Informatika.  
		\item Bapak Martin C. T. Manullang, Ph.D. selaku Dosen Pembimbing atas ide, waktu, tenaga, perhatian, dan masukan yang telah disumbangsihkan kepada penulis.
            \item Bapak Andika Setiawan, S. Kom., M. Cs dan Bapak Eko Dwi Nugroho, S.Kom., M.Cs. selaku Dosen Penguji atas saran dan masukan yang diberikan. 
		\item Kedua Orang Tua dan Adik yang selalu memberikan dukungan dan doa sehingga penulis dapat menyelesaikan tugas akhir ini. Kelembutan, dukungan, dan cinta yang kalian berikan selalu menjadi sumber inspirasi dan kekuatan.
		\item Teman-teman penulis yang membantu selama masa perkuliahan yang tidak bisa disebutkan satu persatu.
	\end{enumerate} \par
	Akhir kata penulis berharap semoga tugas akhir ini dapat memberikan manfaat bagi kita semua. Penulis menyadari bahwa tugas akhir ini tidak luput dari kekurangan dan kelemahan, dan penulis  terbuka untuk menerima saran, kritik, dan masukan yang membangun.
	\vfill
	
\end{justifying}
\clearpage
\end{comment}